\scanpage{76}
Formes quadratiques

(Exposé de Kneser à l’IHES)
\margin{\textcircled{1} \uline{Généralités}}

$A$ anneau commutatif

$E(A)$ anneau des formes quadratiques ``non dégénérées'' [i.e. à discriminant non div de $0$] sur modules projectifs de type fini

$F(A)$ anneau des formes quadratiques ``unimodulaires'' [i.e. à discriminant inversible]

Si $A$ est un corps, $E(A)=F(A)$

\uline{Invariants pour} $F(A)$

1°) Rang \quad $\epsilon:F(A)\longrightarrow\mathbf Z$ \quad hom d’anneaux

2°) Discriminant \quad $d:F(A)\longrightarrow A^*/A^{*2}$ \quad hom de groupes

3°) Classe de Clifford \quad $c:F(A)\longrightarrow{}_2\operatorname{Br}(A)$ \quad application

La définition de $c$ par Kneser se fait en le définissant d’abord pour un module quadratique [en distinguant le cas pair, on prend $\operatorname{Cl}(E)$, et impair, on prend $\operatorname{Cl}^{+}(E)$], puis on \unsure{note} que si $E$ et $E'$ ont \deleted{de rang} discriminant unité, alors $c(E+E')=c(E)+c(E')$, donc on peut définir $c$ au moins sur le sous-groupe $F^2(A)$ noyau de $\epsilon,d$ ----

[N.B. Lorsque $A$ est local de car rés $\neq2$, on peut donner une définition des $w_i$ de Stiefel-Whitney, redonnant $d$ et $c$ -----
\margin{75}
